\documentclass[UTF8]{ctexart}
\usepackage{xeCJK}
\usepackage{geometry}
\geometry{left=2cm,right=2cm,top=2.5cm,bottom=2.5cm}
\setlength{\headheight}{12.7pt}

\usepackage{titlesec}
\titleformat{\section}
  {\normalfont\large\bfseries}
  {\thesection}
  {1em}
  {}
\titlespacing*{\section}
  {0pt}
  {3.5ex plus 1ex minus .2ex}
  {2.3ex plus .2ex}

\usepackage{amsmath}
\usepackage{amssymb}
\usepackage{amsthm}
\usepackage{enumitem}
\setlist[enumerate]{itemsep=0pt,topsep=0pt,parsep=0pt,leftmargin=0pt,itemindent=2.5em}
\setlist[itemize]{itemsep=0pt,topsep=0pt,parsep=0pt,leftmargin=0pt,itemindent=2.5em}
\usepackage{fancyhdr}
\pagestyle{fancy}
\fancyhf{}
\usepackage{graphicx}
\usepackage{hyperref}
\usepackage{indentfirst}
\usepackage{lmodern}

\begin{document}
\setlength{\abovedisplayskip}{1pt}
\setlength{\belowdisplayskip}{1pt}

\section{良序原理}

\hypertarget{sec:定理1.1}{}
\noindent\textbf{定理1.1 (良序原理)}
\vspace{5pt}

任意集合上都有\href{../01 序关系/04 良序关系#nameddest=sec:定义1.1}{良序关系}.

\vspace{7pt}\noindent{\kaishu 证明.}

任取集合$A$, 再取$\mathcal{P}(A)\setminus\{\varnothing\}$上的选择函数$\varphi$,
即对$A$的任意非空子集$a$, 有$\varphi(a)\in a$.

据此超限递归定义$\mathbf{On}$上的类映射$\mathbf{P}$ :
\vspace{3pt}
\[
    \forall\alpha\in\mathbf{On}:\quad\mathbf{P}(\alpha)=\left\{\begin{aligned}
        &\varphi\,(A\setminus\mathbf{P}[\alpha]\,),\quad&&
        \mathbf{P}[\alpha]\text{\,\,是\,\,}A\text{\,\,的真子集\,},\\[-3pt]
        &A,\quad&&\mathbf{P}[\alpha]\text{\,\,不是\,\,}A\text{\,\,的真子集\,}.
    \end{aligned}\right.
\]

\vspace{10pt}\hspace{-1em}\noindent{\kaishu 命题1.}
任给序数$\alpha,\beta$, 如果$\mathbf{P}[\alpha],\mathbf{P}[\beta]$都是$A$的真子集
并且$\mathbf{P}(\alpha)=\mathbf{P}(\beta)$, 那么$\alpha=\beta$.

假设$\beta<\alpha$, 则$\mathbf{P}(\alpha)\in A\setminus\mathbf{P}[\alpha]$,
$\mathbf{P}(\beta)\in\mathbf{P}[\alpha]$, 矛盾.

同理假设$\alpha<\beta$也矛盾. 所以$\alpha=\beta$.

\vspace{7pt}\hspace{-1em}\noindent{\kaishu 命题2.}
存在序数$\alpha$使得$\mathbf{P}[\alpha]$不是$A$的真子集.

假设$\mathbf{P}[\alpha]$都是$A$的真子集, 则$\mathbf{P}(\alpha)$
都属于$A$, 同时由命题1得$\mathbf{P}$是单射, 矛盾.

\vspace{7pt}
把满足命题2的最小序数记作$\mu$, 即:
\begin{enumerate}
    \item $\mathbf{P}[\mu]$不是$A$的真子集,
    \item 对任意的$\beta<\mu$, $\mathbf{P}[\beta]$是$A$的真子集,
    从而$\mathbf{P}(\beta)\in A$.
\end{enumerate}

\vspace{7pt}\hspace{-1em}\noindent{\kaishu 命题3.}
$\mathbf{P}\!\restriction_\mu\,:\mu\to A$是双射.

显然$\mathbf{P}\!\restriction_\mu\,:\mu\to A$, 由命题1得它是单射.
同时$\mathbf{P}[\mu]$不是$A$的真子集, 所以$\mathbf{P}[\mu]=A$.

\vspace{7pt}
最后, $\mathbf{P}\!\restriction_\mu$把严格良序结构$(\mu,\in)$复制到$A$上,
从而$A$上有良序. \hfill $\qedsymbol$

\end{document}
